\section{선별 문제 풀이 I}
\label{sec:symmetric-discriminant-solutions-a}

\subsection*{문제 \thechapter.1: 기본대칭식으로 바꾸기}

\begin{solution}
세 변수에서
\[
  s_1=T_1+T_2+T_3,
  \quad
  s_2=T_1T_2+T_1T_3+T_2T_3,
  \quad
  s_3=T_1T_2T_3
\]
이다.

(a) $s_1^2$을 전개하면
\[
  s_1^2=T_1^2+T_2^2+T_3^2+2s_2
\]
이므로
\[
  T_1^2+T_2^2+T_3^2=s_1^2-2s_2.
\]

(b) 세제곱합의 표준 항등식으로
\[
  T_1^3+T_2^3+T_3^3
  =s_1^3-3s_1s_2+3s_3.
\]
직접 전개하면 $s_1^3$의 혼합항 가운데 $3s_1s_2$를 빼고 $3s_3$을 되돌려 주어야 함을 확인할 수 있다.

(c)
\[
  s_2^2
  =T_1^2T_2^2+T_1^2T_3^2+T_2^2T_3^2
  +2T_1T_2T_3(T_1+T_2+T_3)
\]
이므로
\[
  T_1^2T_2^2+T_1^2T_3^2+T_2^2T_3^2
  =s_2^2-2s_1s_3.
\]

(d)
\[
  s_1s_2
  =\sum_{i\ne j}T_i^2T_j+3T_1T_2T_3
\]
이므로
\[
  \sum_{i\ne j}T_i^2T_j=s_1s_2-3s_3.
\]
\end{solution}

\subsection*{문제 \thechapter.4: 판별식과 중근}

\begin{solution}
(a) $f=X^2-6X+9$에서
\[
  \Disc(f)=(-6)^2-4\cdot9=0.
\]
실제로 $f=(X-3)^2$이다.

(b) $f=X^3-X+1$은 우울형 삼차식이고
\[
  \Disc(f)=-4(-1)^3-27(1)^2=4-27=-23.
\]
판별식이 영이 아니므로 중근이 없다.

(c) $f=X^3-3X+2$에 대해
\[
  \Disc(f)=-4(-3)^3-27(2)^2=108-108=0.
\]
실제로
\[
  f=(X-1)^2(X+2)
\]
이다.

(d) $f=X^4-2$에 $X^m+aX+b$ 공식을 $m=4$, $a=0$, $b=-2$로 적용하면
\[
  \Disc(f)=4^4(-2)^3=-2048.
\]
따라서 중근이 없다. 도함수 $4X^3$과의 최대공약수가 $1$인 것으로도 확인할 수 있다.
\end{solution}

\subsection*{문제 \thechapter.5: 기본 종결식 계산}

\begin{solution}
(a)
\[
  \Res(aX+b,cX+d)
  =\det\begin{pmatrix}a&b\\c&d\end{pmatrix}
  =ad-bc.
\]

(b) $g=X-r$의 유일한 근은 $r$이고 두 다항식의 차수곱이 짝수이므로
\[
  \Res(X^2+pX+q,X-r)
  =r^2+pr+q.
\]
직접 근 공식으로도
\[
  (\alpha_1-r)(\alpha_2-r)=f(r)
\]
를 얻는다.

(c) $X^2-a$의 두 근에서 $X^2-b$의 값은 모두 $a-b$이므로
\[
  \Res(X^2-a,X^2-b)=(a-b)^2.
\]

(d) $X^m-a$의 근들의 곱은 비에타 공식으로
\[
  \prod_{i=1}^m\alpha_i=(-1)^m(-a)=(-1)^{m+1}a.
\]
따라서
\[
  \Res(X^m-a,X)=(-1)^{m+1}a.
\]
\end{solution}

\subsection*{문제 \thechapter.6: 공통근 판정}

\begin{solution}
(a)
\[
  \Res(X^2-2,X^2-3)=(2-3)^2=1\ne0.
\]
따라서 공통근이 없다.

(b) $X^3-2X=X(X^2-2)$이므로 공통인수 $X^2-2$를 갖는다. 따라서 종결식은 $0$이다.

(c)
\[
  X^3-X=X(X-1)(X+1),
  \qquad
  X^2-1=(X-1)(X+1)
\]
이므로 두 근 $\pm1$을 공유하고 종결식은 $0$이다.

(d) $\F_2$에서
\[
  3X^2+1=X^2+1=(X+1)^2.
\]
그러나
\[
  (X^3+X+1)\big|_{X=1}=1+1+1=1
\]
이므로 공통근이 없다. 실제로
\[
  X^3+X+1=X(X^2+1)+1
\]
이어서 최대공약수는 $1$이고 종결식은 $1\in\F_2$이다.
\end{solution}

\subsection*{문제 \thechapter.7: 기약 삼차식의 갈루아군}

\begin{solution}
삼차식은 유리근이 없으면 $\Q$ 위에서 기약이다.

(a) $X^3-X+1$은 $\pm1$을 근으로 갖지 않고
\[
  \Disc=-23
\]
은 제곱이 아니다. 따라서 갈루아군은 $S_3$이다.

(b) $X^3-3X+1$도 $\pm1$을 근으로 갖지 않고
\[
  \Disc=81=9^2.
\]
따라서 갈루아군은 $A_3\cong C_3$이다.

(c) $X^3-2$는 소수 $2$에 대한 아이젠슈타인 판정법으로 기약이다. 판별식은
\[
  -27(-2)^2=-108
\]
이고 $\Q$에서 제곱이 아니므로 갈루아군은 $S_3$이다.

(d) $f=X^3+3X^2-1$은 $\pm1$을 근으로 갖지 않는다. 일반 삼차식 공식에서 $a=3$, $b=0$, $c=-1$을 대입하면
\[
  \Disc(f)=-4(3)^3(-1)-27(-1)^2=108-27=81.
\]
따라서 갈루아군은 $C_3$이다.
\end{solution}

\subsection*{문제 \thechapter.11: 최고 단항식 보조정리}

\begin{solution}
$s_j$의 사전식 최고 단항식은
\[
  T_1T_2\cdots T_j
\]
이고 그 계수는 $1$이다. 따라서
\[
  M=s_1^{\lambda_1-\lambda_2}
  s_2^{\lambda_2-\lambda_3}\cdots s_n^{\lambda_n}
\]
의 최고 단항식은 각 최고 단항식을 곱한
\[
  \prod_{j=1}^n
  (T_1\cdots T_j)^{\lambda_j-\lambda_{j+1}},
  \qquad \lambda_{n+1}=0
\]
이다. 여기서 $T_i$의 지수는
\[
  \sum_{j=i}^n(\lambda_j-\lambda_{j+1})
  =\lambda_i
\]
이므로 최고 단항식은 $T_1^{\lambda_1}\cdots T_n^{\lambda_n}$이다. 계수는 $1$들의 곱이므로 $1$이다.
\end{solution}

\subsection*{문제 \thechapter.14: 반복나눗셈 자유기저}

\begin{solution}
$h\in A[X]$가 일계수 $d$차이므로 임의의 $f\in A[X]$에 대해 나눗셈 알고리즘을 반복할 수 있다.
\[
  f=f_1h+r_0,
  \quad
  f_1=f_2h+r_1,
  \quad\dots,
  \qquad \deg r_i<d.
\]
몫의 차수는 매 단계 감소하므로 유한 단계 뒤에 끝나며
\[
  f=r_0+r_1h+\cdots+r_mh^m
\]
를 얻는다. 각 $r_j$를
\[
  r_j=a_{0j}+a_{1j}X+\cdots+a_{d-1,j}X^{d-1}
\]
로 쓰면
\[
  f=F_0(h)+F_1(h)X+\cdots+F_{d-1}(h)X^{d-1}
\]
인 $F_i\in A[Y]$를 얻는다.

유일성을 보자. 만약
\[
  0=r_0+r_1h+\cdots+r_mh^m,
  \qquad \deg r_j<d
\]
라면 $h$에 대한 나눗셈의 나머지가 $r_0$이므로 $r_0=0$이다. $h$를 묶고 같은 논증을 반복하면 모든 $r_j=0$이다. 따라서
\[
  A[X]=\bigoplus_{i=0}^{d-1}A[h]X^i.
\]
\end{solution}

\subsection*{문제 \thechapter.16: 종결식 베주 항등식}

\begin{solution}
실베스터 사상
\[
  \Phi:R[X]_{<n}\oplus R[X]_{<m}	o R[X]_{<m+n},
  \qquad (p,q)\mapsto pf+qg
\]
의 행렬을 $M$이라 하자. $M$은 실베스터 행렬의 전치이므로
\[
  \det M=\Res(f,g).
\]
수반행렬 항등식
\[
  M\operatorname{adj}(M)=(\det M)I
\]
에서 공역의 상수다항식 $1$에 대응하는 기저벡터 $e$를 택한다. 정의역 벡터
\[
  v=\operatorname{adj}(M)e
\]
를 두 성분 $(p,q)$로 해석하면
\[
  \Phi(p,q)=Mv=(\det M)e=\Res(f,g).
\]
정의역 자체가 $R[X]_{<n}\oplus R[X]_{<m}$이므로
\[
  \deg p<n,
  \qquad
  \deg q<m
\]
도 자동으로 성립한다.
\end{solution}
